\section{Formalization Design and Reuse}
\label{sec:engineering}

\paragraph{Shape of the development.}
Fifteen Lean~4 modules, about 4{,}300 lines, over a pinned
\textsf{mathlib} (both at v4.31.0): the core valuation theory and
GMT bridge, the DS layer (2-monotonicity, inclusion--exclusion and
the $\infty$-monotone tower, the two conditionings), the
representation suite (finite, obstruction, Scott, general
decomposition), the Cox layer (\texttt{CoxModel},
sum-irreducibility, the corrected theorem), the Acz\'el development,
the two guards (halting chain, Sierpi\'nski model), the product rigidity
theorem on chain products, the countable-mixture composition laws,
and the ultrafilter couplings that separate the iteration orders.
The build is clean, with no
\texttt{sorry} and no warnings, and every theorem named in this
paper reports exactly
\texttt{[propext, Classical.choice, Quot.sound]} under
\texttt{\#print axioms}.

\paragraph{Meta-logic against object logic.}
Lean's logic is classical, while our \emph{object} of study is the
frame.  This division of labor is deliberate and, we found, the
pragmatic sweet spot.  Constructivity claims attach to the algebra
(no theorem assumes $a \sqcup a^{\mathsf{c}} = \top$), while the
meta-theory may reason classically \emph{about} it.  The
counterexample valuations of Sections~\ref{sec:representation}
and~\ref{sec:cox} use classical case splits legitimately, since a
counterexample is a classical fact about the theory.  Where the
distinction has content we track it: the non-vacuity witness for
\texttt{CoxModel} was re-based from \texttt{Prop} (which imports
excluded middle through the back door) to the chain $\ENNReal$,
where the needed case split is a constructor check.

\paragraph{What the checker contributed.}
Beyond the usual assurance, five findings in this paper originate in
failed proofs.  The naive \texttt{constructive\_cox} statement was
unsatisfiable (Section~\ref{sec:cox}).  Unrestricted strict
monotonicity for the conjunction functional was vacuous
(Section~\ref{sec:cox}).  The regular-elements guess for the
classical fragment was false (Section~\ref{sec:valuation}).  Our
first Acz\'el hypothesis bundle omitted order-preservation, which
the additivity proof needs (Appendix~\ref{app:aczel}).  And our
expected counterexample to product uniqueness was refuted by the
rigidity that became Section~\ref{sec:conclusion}'s theorem.  We
regard these as the paper's best evidence that this subject
specifically repays mechanization.

\paragraph{\textsf{mathlib} surface.}
The development consumes order theory
(\texttt{Order.Frame}, Heyting algebras, prime ideal separation),
\texttt{ENNReal} arithmetic, measure theory for the bridge,
\texttt{PMF}, and the computability library
(\texttt{Nat.Partrec.Code}, \texttt{ComputablePred}).  Candidates
for upstreaming, in decreasing generality: the Acz\'el/H\"older
generator construction of Appendix~\ref{app:aczel}
(\textsf{mathlib} has no associativity-representation theorem, no
monotone Cauchy functional equation, and no t-norm theory),
valuations on frames with the mass/representation API, and the
inclusion--exclusion identity for modular functions on distributive
lattices (Theorem~\ref{thm:incl-excl}), which is independent of
probability.  $\ENNReal$'s truncated subtraction deserves a note.
It forces additive normal forms, with every identity stated so that
no subtraction occurs.  This reads as a constraint but functioned as
a correctness discipline: several informal manipulations of the DS
literature do not survive it unchanged.

\paragraph{Artifact.}
The artifact (anonymized supplement) builds with a two-command
recipe on the pinned toolchain and re-checks every theorem,
including the axiom audit.
